\section{哀伤、丧亲与心理支持资源}

丧失至亲是人生最沉重的经历之一，而它常常同时冲击着身体、情绪、关系与性。这一节把\textbf{哀伤的正常过程、何时需要专业帮助，以及可用的支持资源}整理在一起，供经历者与其身边人取用。

\subsection{先理解：哀伤是正常的，没有"标准进度表"}

\begin{itemize}
	\item \textbf{哀伤不是病}：它是面对丧失的自然反应，每个人的节奏、方式与时长都不同，不应用"你怎么还没走出来"来催促。
	\item \textbf{"哀伤五阶段"不是必经流程}：否认、愤怒、讨价还价、抑郁、接受这一模型是\textbf{描述性的}，并非人人按顺序经历；把它当"考纲"反而会增加压力。
	\item \textbf{哀伤可能反复}：纪念日、节日、熟悉的地点都可能唤起强烈的思念与痛苦，这属正常波动，不代表"退步"。
\end{itemize}

\subsection{哀伤期的性与亲密}

\begin{itemize}
	\item \textbf{反应两极都正常}：有人在丧亲后性欲明显下降，也有人因孤独、需要慰藉而性需求增强——\textbf{两种都是正常的心理生理反应}，不必自责，更不应被他人污名化。
	\item \textbf{"对逝者不忠"是误解}：寻求亲密与慰藉，与对逝者的爱与怀念并不矛盾。
	\item \textbf{若以性行为作为唯一的情绪出口}并逐渐失控、带来损害，则应寻求专业帮助（参见"女性性欲亢进"相关内容）。
\end{itemize}

\subsection{可用的支持资源}

\begin{itemize}
	\item \textbf{全国统一心理援助热线 12356}：24 小时、免费、匿名，适合情绪困扰、丧亲哀伤、亲密关系危机等（详见"心理危机与情绪支持"一节）。
	\item \textbf{综合医院心理科 / 精神科}：评估是否合并抑郁、焦虑、失眠或"延伸哀伤障碍"（详见下文）。
	\item \textbf{安宁疗护 / 临终关怀机构}：通常提供\textbf{遗属随访与哀伤支持}，可向曾提供临终照护的机构咨询。
	\item \textbf{专业协会与公益组织}：各地心理卫生协会、宁养院、哀伤辅导机构等；部分提供个体咨询与团体辅导。
	\item \textbf{互助支持团体}：丧偶、失独、丧子父母等同伴互助团体，能让经历者"不必解释就被理解"，常是重要的慰藉来源。
\end{itemize}

\subsection{自助与陪伴：怎么做才真正有帮助}

\begin{itemize}
	\item \textbf{允许情绪存在}：想哭就哭，想说就说；压抑与"假装坚强"往往让痛苦更久。
	\item \textbf{重建基本节律}：规律作息、适度运动、正常饮食——身体稳定是情绪恢复的基础；尽量避免以酒精或过量行为麻痹自己。
	\item \textbf{给自己时间，也保持连接}：与可信任的人保持联系，参加支持团体，不必独自硬扛。
	\item \textbf{给身边人的建议}：\textbf{陪伴大于劝导}。"你要坚强""时间会治愈""想开点"等安慰往往无效甚至伤人；一句"我在这里，陪你"和持续的实际帮助更有力量。特别在\textbf{纪念日、节假日与头一年}等高风险期主动关心。
\end{itemize}

\subsection{何时需要专业帮助}

\begin{itemize}
	\item \textbf{延伸哀伤障碍（Prolonged Grief Disorder）}：成人哀伤持续约 \textbf{12 个月以上}（儿童青少年约 6 个月），仍强烈影响日常生活功能，可考虑为"延伸哀伤"，需要专业干预（如哀伤聚焦的认知行为治疗、人际治疗）。
	\item \textbf{合并抑郁 / 创伤}：持续情绪低落、失眠、进食障碍、对一切失去兴趣，或出现创伤性再体验，应及时就诊心理科/精神科。
	\item \textbf{红旗信号——立即求助}：出现\textbf{自杀念头、自伤行为或无法自理}时，请\textbf{立即拨打 12356 或 120/110}，并请身边人陪伴，不要独处。
\end{itemize}

\begin{tcolorbox}[colback=purple!5!white,colframe=purple!55!black,title=贴心小叮咛]
哀伤不是要被"克服"的敌人，而是要被\textbf{陪伴着走过去}的一段路。它没有统一的时间表，也不该被催促。若你正身处其中，请允许自己慢一点，也别忘了伸手求助——热线、心理科、安宁疗护团队与同伴团体，都是可以依靠的资源。你不需要一个人扛下全部。
\end{tcolorbox}
